
\section{Enfoque, tipo y diseño metodológico}

\subsection{Enfoque cuantitativo}

Se recomienda el enfoque cuantitativo porque la pregunta de investigación busca describir frecuencias, categorías de herramientas digitales y valoraciones recogidas mediante respuestas estructuradas. Este enfoque permite codificar la información, resumirla mediante tablas y porcentajes y comparar descriptivamente grupos cuando la cantidad y la calidad de los datos lo permitan. La elección del enfoque debe guardar coherencia con el problema, los objetivos y los procedimientos previstos para recolectar y analizar la información \parencite{feriaavila2019dimension,guerrero2015metodologia}.

La recomendación no significa que sea el único enfoque posible. Si el propósito central fuera comprender en profundidad las experiencias personales y los significados atribuidos al uso de estas herramientas, sería pertinente considerar un diseño cualitativo. Si se combinaran mediciones estructuradas y entrevistas, podría plantearse un enfoque mixto.

\subsection{Tipo y alcance}

La investigación será aplicada en su orientación práctica y descriptiva en su alcance. Se considera descriptiva porque pretende caracterizar las prácticas de uso de herramientas digitales con fines académicos, sin probar una intervención ni establecer relaciones causales. Puede incorporar comparaciones exploratorias entre subgrupos; sin embargo, estas deberán presentarse como análisis descriptivos y no como evidencia de causalidad.

\subsection{Diseño no experimental y transversal}

El diseño será no experimental, puesto que no se manipularán deliberadamente las variables ni se asignará a los participantes a tratamientos. Será transversal porque la información se recogerá una sola vez durante un periodo previamente definido. Esta elección resulta viable para una caracterización inicial, aunque no permite observar cambios a lo largo del tiempo ni demostrar que una herramienta produzca mejoras en el rendimiento académico.

\subsection{Método y fases}

El procedimiento metodológico se organizará en seis fases:

\begin{enumerate}
    \item \textbf{Delimitación:} seleccionar la institución o las instituciones participantes, el periodo académico y la población accesible; además, obtener las autorizaciones necesarias.

    \item \textbf{Diseño del instrumento:} elaborar un cuestionario alineado con las variables de estudio y someterlo a revisión para evaluar la claridad y pertinencia de las preguntas.

    \item \textbf{Prueba piloto:} aplicar el cuestionario a un grupo pequeño con características similares a las de la población objetivo, procurando que sus respuestas no se incluyan en el análisis principal cuando sea posible. Ajustar las preguntas ambiguas a partir de los resultados de esta prueba.

    \item \textbf{Recolección de datos:} invitar a los estudiantes a participar voluntariamente, informarles sobre el propósito del estudio y obtener su consentimiento informado.

    \item \textbf{Depuración y análisis:} revisar los registros incompletos, codificar las respuestas y calcular las estadísticas descriptivas correspondientes. La definición de la muestra y su alcance deberán documentarse de manera explícita, especialmente si se recurre a participantes voluntarios \parencite{argibay2009muestra}.

    \item \textbf{Elaboración del informe:} presentar los resultados obtenidos, las limitaciones del estudio y las recomendaciones que se deriven de la evidencia recopilada.
\end{enumerate}

\subsection{Razones para elegir esta metodología}

La metodología propuesta responde directamente a la pregunta de investigación, emplea variables observables y utiliza un instrumento de aplicación relativamente breve. Asimismo, facilita resumir las respuestas de distintos estudiantes y puede ejecutarse con recursos académicos ordinarios.

Su principal limitación es que las respuestas serán autoinformadas y podrían verse afectadas por errores de memoria, interpretaciones diferentes de las opciones o deseabilidad social. La selección muestral también condicionará el alcance de las conclusiones, por lo que deberá distinguirse entre los resultados obtenidos en la muestra y las posibles inferencias hacia una población más amplia \parencite{argibay2009muestra}.

\subsection{Coherencia metodológica}

La correspondencia entre los componentes principales del estudio se establece de la siguiente manera:


\begin{callout}
    \textbf{Coherencia metodológica}
    \smallskip
    \centering
    \textbf{Pregunta descriptiva}
    \(\rightarrow\)
    \textbf{Encuesta estructurada}
    \(\rightarrow\)
    \textbf{Variables observables}
    \(\rightarrow\)
    \textbf{Estadística descriptiva}
    \(\rightarrow\)
    \textbf{Conclusiones delimitadas}
\end{callout}
